\documentclass[10pt]{article}
\usepackage{amsmath,amssymb,amsthm}
\usepackage[margin=1in]{geometry}

\newtheorem{theorem}{Theorem}
\newtheorem{proposition}{Proposition}
\newtheorem{lemma}{Lemma}
\newtheorem{example}{Example}
\newcommand{\R}{\mathbb R}

\title{Exact Bellman Values with Euclidean Norm Costs and Conic Control Bounds}
\author{Anders Hansson}
\date{}

\begin{document}
\maketitle

\begin{abstract}
An infinite-horizon control problem is solved on a second-order cone. The dynamics are linear, admissible state-control triples satisfy finitely many conic inequalities, and the stage cost contains a Euclidean norm of the vector state component and control together. An explicit parameter condition gives a value function equal to a linear term plus a Euclidean norm, with an optimal feedback that leaves a nonzero state at every finite time. A two-variable conic program computes the value coefficients. The condition on the control bound is sharp for this value-function form: below the threshold, the value differs from that form at boundary states. The results provide an exact example beyond linear values on the original state cone, while identifying where this simple formula ceases to apply.
\end{abstract}

\section{Introduction}

Pates and Rantzer \cite{pates-rantzer} give conditions under which a control problem on a positive cone has a linear value function computed by convex optimization. Bellman inequalities also give lower bounds when the value is sought among functions parameterized by finitely many coefficients \cite{de-farias-van-roy}. For an exact extension, one such function must equal the value, and an admissible policy must attain its lower bound.

This paper considers a value of the form $a t+b\|z\|_2$ on states $(t,z)$ satisfying $t\geq\|z\|_2$. The norm term makes the value nonlinear on this cone when $b>0$ and $z\in\R^n$ with $n\geq2$. The stage cost uses one Euclidean norm of the stacked state and control terms. This choice permits an interior minimizing control. Separate state and control norms can instead give a scalar objective whose minimum lies at an endpoint, producing no control or immediate cancellation.

The conic control bound matters even though it need not bind at the optimal policy. A threshold on that bound determines exactly when the displayed value formula is valid for the problem studied here. Below the threshold, the value agrees with the formula at sufficiently interior states and is strictly larger at boundary states. This gives both an exact conic formulation and a test of its scope.

\section{Control Problem and Bellman Lower Bound}

Fix an integer $n\geq2$ and define the second-order cone
\begin{equation}
\mathcal K=\{(t,z)\in\R\times\R^n:t\geq\|z\|_2\}.
\label{eq:state-cone}
\end{equation}
Let $P\in\R^{n\times n}$ satisfy $P^TP=I$. Fix
\begin{equation}
0<\alpha<1,\qquad 0<\gamma\leq1,\qquad
c>0,\quad r>0,\quad s>0,\quad\kappa>0.
\label{eq:data}
\end{equation}
The admissible state-control triples form the cone
\begin{equation}
\mathcal G_\kappa=\{(t,z,u):(t,z)\in\mathcal K,\quad
\|u\|_2\leq\kappa t,\quad
\|\alpha z+u\|_2\leq\alpha t\}.
\label{eq:control-cone}
\end{equation}
For $(t_k,z_k,u_k)\in\mathcal G_\kappa$, the successor and stage cost are
\begin{equation}
t_{k+1}=\alpha t_k,\qquad
z_{k+1}=P(\alpha z_k+u_k),\qquad
\ell(t_k,z_k,u_k)=c t_k+
\sqrt{r^2\|z_k\|_2^2+s^2\|u_k\|_2^2}.
\label{eq:dynamics-cost}
\end{equation}
The value at $(t_0,z_0)\in\mathcal K$ is
\begin{equation}
J(t_0,z_0)=\inf_{\{u_k\}_{k\geq0}}
\sum_{k=0}^{\infty}\gamma^k\ell(t_k,z_k,u_k),
\quad (t_k,z_k,u_k)\in\mathcal G_\kappa\text{ for every }k.
\label{eq:value}
\end{equation}
The control $u=0$ is feasible at every state in \eqref{eq:state-cone}. The last inequality in \eqref{eq:control-cone} makes the successor in \eqref{eq:dynamics-cost} belong to \eqref{eq:state-cone}. The zero-control trajectory has $t_k=\alpha^kt_0$ and $\|z_k\|_2=\alpha^k\|z_0\|_2$, so its cost is $(ct_0+r\|z_0\|_2)/(1-\gamma\alpha)<\infty$. Each condition in \eqref{eq:control-cone} is a second-order cone constraint. The cone is closed and convex. It is pointed because membership of both $(t,z,u)$ and $-(t,z,u)$ forces $t=z=u=0$. It has nonempty interior because all inequalities are strict at $(t,z,u)=(1,0,0)$.

For $a,b\geq0$, define the candidate function
\begin{equation}
V_{a,b}(t,z)=a t+b\|z\|_2,\qquad (t,z)\in\mathcal K.
\label{eq:candidate}
\end{equation}
For any function $V$ on \eqref{eq:state-cone}, define its Bellman image by
\begin{equation}
(TV)(t,z)=\inf_{u:(t,z,u)\in\mathcal G_\kappa}
\bigl(\ell(t,z,u)+\gamma V(\alpha t,P(\alpha z+u))\bigr).
\label{eq:bellman-operator}
\end{equation}
We call $V$ a Bellman subsolution if $V\leq TV$ on $\mathcal K$.

\begin{lemma}
\label{lem:lower-bound}
Let $a,b\geq0$. If $V_{a,b}\leq TV_{a,b}$ on $\mathcal K$, then $V_{a,b}\leq J$ on $\mathcal K$.
\end{lemma}

\begin{proof}
Fix an admissible trajectory from $(t_0,z_0)$. By \eqref{eq:bellman-operator}, the hypothesis gives
\begin{equation}
V_{a,b}(t_k,z_k)\leq
\ell(t_k,z_k,u_k)+\gamma V_{a,b}(t_{k+1},z_{k+1})
\label{eq:one-step}
\end{equation}
for every $k$. Multiply \eqref{eq:one-step} by $\gamma^k$ and sum from $k=0$ to $N$ to obtain
\begin{equation}
V_{a,b}(t_0,z_0)\leq
\sum_{k=0}^{N}\gamma^k\ell(t_k,z_k,u_k)
+\gamma^{N+1}V_{a,b}(t_{N+1},z_{N+1}).
\label{eq:telescope}
\end{equation}
Equations \eqref{eq:state-cone} and \eqref{eq:dynamics-cost} give
\begin{equation}
0\leq\gamma^{N+1}V_{a,b}(t_{N+1},z_{N+1})
\leq(a+b)(\gamma\alpha)^{N+1}t_0\longrightarrow0,
\label{eq:tail-bound}
\end{equation}
because $\gamma\alpha<1$ by \eqref{eq:data}. Taking $N\to\infty$ in \eqref{eq:telescope} and then the infimum over admissible trajectories in \eqref{eq:value} proves the claim.
\end{proof}

\section{Control Minimization}

For $b\geq0$ and $0\leq\theta\leq\alpha$, define
\begin{equation}
g_b(\theta)=\sqrt{r^2+s^2\theta^2}+\gamma b(\alpha-\theta).
\label{eq:scalar-objective}
\end{equation}

\begin{proposition}
\label{prop:reduction}
For $a,b\geq0$ and $(t,z)\in\mathcal K$ with $z\ne0$,
\begin{equation}
(TV_{a,b})(t,z)
=(c+\gamma\alpha a)t+
\|z\|_2\min_{0\leq\theta\leq
\min\{\alpha,\kappa t/\|z\|_2\}}g_b(\theta).
\label{eq:reduced-bellman}
\end{equation}
At $z=0$, $(TV_{a,b})(t,0)=(c+\gamma\alpha a)t$.
\end{proposition}

\begin{proof}
Orthogonality of $P$ in \eqref{eq:data} gives
\begin{equation}
\|P(\alpha z+u)\|_2=\|\alpha z+u\|_2.
\label{eq:orthogonal-norm}
\end{equation}
Fix $z\ne0$ and set $\rho=\|z\|_2$. If a feasible $u$ has $\|u\|_2>\alpha\rho$, replace it by $-\alpha z$. This replacement has a smaller norm and satisfies $\|-\alpha z\|_2=\alpha\rho<\|u\|_2\leq\kappa t$; its successor has norm zero. Hence it is feasible in \eqref{eq:control-cone} and reduces both norm terms in \eqref{eq:bellman-operator}.

Now let a feasible $u$ have $v=\|u\|_2\leq\alpha\rho$ and replace it by $-v z/\rho$. The replacement has norm $v$, so its stage cost and control bound are unchanged. The reverse triangle inequality gives
\begin{equation}
\|\alpha z+u\|_2\geq\alpha\rho-v
=\|\alpha z-vz/\rho\|_2.
\label{eq:radial-bound}
\end{equation}
Thus the replacement satisfies the successor constraint in \eqref{eq:control-cone} and does not increase the continuation term in \eqref{eq:bellman-operator}. Every minimizing control can therefore be chosen as $u=-\theta z$, with $0\leq\theta\leq\min\{\alpha,\kappa t/\rho\}$. Such controls satisfy the successor constraint because $(\alpha-\theta)\rho\leq\alpha t$ by \eqref{eq:state-cone}. Substitution in \eqref{eq:bellman-operator} gives \eqref{eq:reduced-bellman}.

If $z=0$, then $u=0$ is feasible and minimizes $s\|u\|_2+\gamma b\|u\|_2$. The remaining terms in \eqref{eq:bellman-operator} give the stated value.
\end{proof}

\section{Exact Value and a Finite Conic Program}

Define the constants
\begin{equation}
d=1-\gamma\alpha,\qquad
L=\sqrt{d^2+(\gamma r/s)^2},\qquad
a_*=\frac{c}{d},\qquad
b_*=\frac{r}{L},\qquad
\eta=\frac{\gamma r^2}{s^2d}.
\label{eq:constants}
\end{equation}
All denominators in \eqref{eq:constants} are positive by \eqref{eq:data}.

\begin{theorem}
\label{thm:exact}
Suppose $\eta<\alpha$ and $\kappa\geq\eta$. Then
\begin{equation}
J(t,z)=a_*t+b_*\|z\|_2,
\qquad (t,z)\in\mathcal K.
\label{eq:exact-value}
\end{equation}
The feedback $u=-\eta z$ is optimal. For $z_0\ne0$, its vector state is nonzero at every finite time.
\end{theorem}

\begin{proof}
For $0\leq\theta\leq\alpha$, differentiation of \eqref{eq:scalar-objective} gives
\begin{equation}
g_{b_*}'(\theta)=
\frac{s^2\theta}{\sqrt{r^2+s^2\theta^2}}-\gamma b_* ,
\qquad
g_{b_*}''(\theta)=
\frac{r^2s^2}{(r^2+s^2\theta^2)^{3/2}}>0.
\label{eq:scalar-derivatives}
\end{equation}
The definitions in \eqref{eq:constants} give $\gamma\eta=(\gamma r/s)^2/d$ and hence
\begin{equation}
L^2=d(d+\gamma\eta),\qquad
\sqrt{r^2+s^2\eta^2}=b_*(d+\gamma\eta),\qquad
\frac{s^2\eta}{\sqrt{r^2+s^2\eta^2}}=\gamma b_*.
\label{eq:coefficient-identities}
\end{equation}
Equations \eqref{eq:scalar-derivatives} and \eqref{eq:coefficient-identities} show that $\eta$ is the unique minimizer of $g_{b_*}$ on $[0,\alpha]$. They also give
\begin{equation}
g_{b_*}(\eta)
=b_*(d+\gamma\eta)+\gamma b_*(\alpha-\eta)
=b_*.
\label{eq:scalar-fixed-point}
\end{equation}

For $z\ne0$, \eqref{eq:state-cone} and $\kappa\geq\eta$ imply $\kappa t/\|z\|_2\geq\eta$. The feedback $u=-\eta z$ also satisfies
\begin{equation}
\|u\|_2=\eta\|z\|_2\leq\kappa t,
\qquad
\|\alpha z+u\|_2=(\alpha-\eta)\|z\|_2\leq\alpha t.
\label{eq:feedback-feasibility}
\end{equation}
By Proposition \ref{prop:reduction}, \eqref{eq:scalar-fixed-point}, and $c+\gamma\alpha a_*=a_*$ from \eqref{eq:constants},
\begin{equation}
TV_{a_*,b_*}=V_{a_*,b_*}\quad\text{on }\mathcal K.
\label{eq:bellman-fixed-point}
\end{equation}
Lemma \ref{lem:lower-bound} therefore gives $V_{a_*,b_*}\leq J$.

Under $u_k=-\eta z_k$, equations \eqref{eq:dynamics-cost} and \eqref{eq:feedback-feasibility} give
\begin{equation}
t_k=\alpha^k t_0,
\qquad
z_k=(\alpha-\eta)^kP^kz_0,
\qquad
\|u_k\|_2=\eta(\alpha-\eta)^k\|z_0\|_2.
\label{eq:feedback-trajectory}
\end{equation}
Since $0<\alpha-\eta<\alpha<1$, \eqref{eq:feedback-trajectory} has $z_k\ne0$ at every finite time when $z_0\ne0$. Its cost is
\begin{equation}
\begin{aligned}
\sum_{k=0}^{\infty}\gamma^k\ell(t_k,z_k,u_k)
&=\frac{c}{1-\gamma\alpha}t_0
+\frac{\sqrt{r^2+s^2\eta^2}}
{1-\gamma(\alpha-\eta)}\|z_0\|_2\\
&=a_*t_0+b_*\|z_0\|_2,
\end{aligned}
\label{eq:feedback-cost}
\end{equation}
where the last equality uses \eqref{eq:constants} and \eqref{eq:coefficient-identities}. Thus $J\leq V_{a_*,b_*}$. Together with the lower bound, this proves \eqref{eq:exact-value} and optimality of the feedback.
\end{proof}

The value coefficients in \eqref{eq:exact-value} can also be obtained without solving the scalar fixed-point equation in \eqref{eq:scalar-fixed-point}. Fix $m_t,m_z>0$ and consider the finite conic program
\begin{equation}
\begin{array}{ll}
\text{maximize} & m_t a+m_z b\\
\text{subject to} & a\geq0,\quad b\geq0,\quad da\leq c,\\
&\displaystyle\left\|
\begin{bmatrix}db/r\\ \gamma b/s\end{bmatrix}
\right\|_2\leq1.
\end{array}
\label{eq:conic-program}
\end{equation}

\begin{proposition}
\label{prop:conic-program}
Every feasible pair in \eqref{eq:conic-program} gives a Bellman subsolution $V_{a,b}$. If $\eta<\alpha$ and $\kappa\geq\eta$, the unique optimizer of \eqref{eq:conic-program} is $(a_*,b_*)$ and its candidate equals $J$.
\end{proposition}

\begin{proof}
The last constraint in \eqref{eq:conic-program} and the Cauchy--Schwarz inequality give, for every $z,u$,
\begin{equation}
\begin{aligned}
db\|z\|_2+\gamma b\|u\|_2
&=\begin{bmatrix}db/r&\gamma b/s\end{bmatrix}
  \begin{bmatrix}r\|z\|_2\\s\|u\|_2\end{bmatrix}\\
&\leq\sqrt{r^2\|z\|_2^2+s^2\|u\|_2^2}.
\end{aligned}
\label{eq:conic-certificate}
\end{equation}
The reverse triangle inequality gives
\begin{equation}
\|\alpha z+u\|_2\geq\alpha\|z\|_2-\|u\|_2.
\label{eq:successor-bound}
\end{equation}
Using \eqref{eq:dynamics-cost}, \eqref{eq:candidate}, \eqref{eq:conic-certificate}, and \eqref{eq:successor-bound}, the Bellman residual satisfies
\begin{equation}
\begin{aligned}
\ell(t,z,u)+\gamma V_{a,b}(\alpha t,P(\alpha z+u))-V_{a,b}(t,z)
&\geq(c-da)t+db\|z\|_2+\gamma b\|u\|_2\\
&\quad+\gamma b(\alpha\|z\|_2-\|u\|_2)-b\|z\|_2\\
&=(c-da)t\geq0.
\end{aligned}
\label{eq:certified-residual}
\end{equation}
The last inequality uses $da\leq c$ in \eqref{eq:conic-program} and $t\geq0$ in \eqref{eq:state-cone}. Thus every feasible pair satisfies $V_{a,b}\leq TV_{a,b}$. Lemma \ref{lem:lower-bound} gives $V_{a,b}\leq J$.

Because $d,r,s>0$, the constraints in \eqref{eq:conic-program} are equivalent to
\begin{equation}
0\leq a\leq\frac{c}{d}=a_*,
\qquad
0\leq b\leq\frac{r}{\sqrt{d^2+(\gamma r/s)^2}}=b_*.
\label{eq:program-bounds}
\end{equation}
Both objective weights are positive, so \eqref{eq:program-bounds} makes $(a_*,b_*)$ the unique optimizer. Under the hypotheses of Theorem \ref{thm:exact}, \eqref{eq:exact-value} identifies its candidate with $J$.
\end{proof}

Inequality \eqref{eq:conic-certificate} is sufficient for the Bellman inequality. The proposition does not assert that it describes every Bellman subsolution of the form \eqref{eq:candidate}. Exactness follows because the feasible pairs in \eqref{eq:conic-program} include the value coefficients under the stated condition on $\kappa$.

\section{Sharpness of the Control Bound}

The candidate $V_{a_*,b_*}$ remains a lower bound when $\kappa<\eta$, since the proof of Proposition \ref{prop:conic-program} does not use $\kappa\geq\eta$. The bound is then strict at states on the boundary of \eqref{eq:state-cone}.

\begin{theorem}
\label{thm:threshold}
Suppose $0<\kappa<\eta<\alpha$. For $z\ne0$ and $t\geq\eta\|z\|_2/\kappa$,
\begin{equation}
J(t,z)=a_*t+b_*\|z\|_2.
\label{eq:interior-value}
\end{equation}
For $\|z\|_2=t>0$,
\begin{equation}
J(t,z)\geq a_*t+g_{b_*}(\kappa)t
>a_*t+b_*t.
\label{eq:boundary-gap}
\end{equation}
No function of the form $a t+b\|z\|_2$ equals $J$ on all of $\mathcal K$.
\end{theorem}

\begin{proof}
Proposition \ref{prop:conic-program} and Lemma \ref{lem:lower-bound} give $V_{a_*,b_*}\leq J$ for every $\kappa>0$. If $t\geq\eta\|z\|_2/\kappa$, then $u=-\eta z$ satisfies the control bound in \eqref{eq:control-cone}. By \eqref{eq:feedback-trajectory},
\begin{equation}
\frac{t_k}{\|z_k\|_2}
=\left(\frac{\alpha}{\alpha-\eta}\right)^k
\frac{t}{\|z\|_2}
\geq\frac{\eta}{\kappa},
\label{eq:ratio-growth}
\end{equation}
so the feedback remains admissible. Equation \eqref{eq:feedback-cost} then gives $J\leq V_{a_*,b_*}$ at these states, proving \eqref{eq:interior-value}.

At a boundary state with $t=\|z\|_2>0$, Proposition \ref{prop:reduction} gives
\begin{equation}
(TV_{a_*,b_*})(t,z)
=a_*t+t\min_{0\leq\theta\leq\kappa}g_{b_*}(\theta)
=a_*t+t g_{b_*}(\kappa).
\label{eq:boundary-bellman}
\end{equation}
The second equality uses $\kappa<\eta<\alpha$ and the strict convexity and unique minimizer established in \eqref{eq:scalar-derivatives}--\eqref{eq:scalar-fixed-point}. Hence $g_{b_*}(\kappa)>g_{b_*}(\eta)=b_*$. For every admissible first control, the remaining cost is at least $J$ at its successor and thus at least $V_{a_*,b_*}$ by Proposition \ref{prop:conic-program}. Taking the infimum over first controls gives $J(t,z)\geq(TV_{a_*,b_*})(t,z)$. Combining this with \eqref{eq:boundary-bellman} proves \eqref{eq:boundary-gap}.

At $z=0$, every stage cost in \eqref{eq:dynamics-cost} is at least $c t_k$. The control $u_k=0$ attains $c t_k$ at every time, so \eqref{eq:value} and $t_k=\alpha^k t$ give $J(t,0)=a_*t$. Any function $a t+b\|z\|_2$ equal to $J$ must therefore have $a=a_*$. Equality on the interior states in \eqref{eq:interior-value} then forces $b=b_*$. This pair contradicts \eqref{eq:boundary-gap}, proving the last claim.
\end{proof}

\section{Numerical Instance}

\begin{example}
Let $n=2$, $P=\begin{bmatrix}0&-1\\1&0\end{bmatrix}$, and choose
\begin{equation}
\gamma=1,\qquad \alpha=\frac12,\qquad
c=1,\qquad r=\frac14,\qquad s=1.
\label{eq:example-data}
\end{equation}
Equation \eqref{eq:constants} gives $d=1/2$, $\eta=1/8$, $a_*=2$, and $b_*=1/\sqrt5$. For every $\kappa\geq1/8$, Theorem \ref{thm:exact} gives
\begin{equation}
J(t,z)=2t+\frac{1}{\sqrt5}\|z\|_2,
\qquad u=-\frac18z.
\label{eq:example-value}
\end{equation}
The trajectory has $t_k=2^{-k}t_0$ and $z_k=(3/8)^kP^kz_0$ by \eqref{eq:feedback-trajectory}, so both components contract while the vector state rotates and remains nonzero at every finite time when $z_0\ne0$.

If $\kappa=1/16$, Theorem \ref{thm:threshold} gives equality with \eqref{eq:example-value} when $t\geq2\|z\|_2$. At a boundary state $t=\|z\|_2>0$, \eqref{eq:boundary-gap} gives the explicit positive gap
\begin{equation}
J(t,z)-\left(2t+\frac{t}{\sqrt5}\right)
\geq\frac{t}{16}\left(\sqrt{17}-\frac9{\sqrt5}\right)>0.
\label{eq:example-gap}
\end{equation}
The strict inequality in \eqref{eq:example-gap} follows from $17>81/5$.
\end{example}

\section{Discussion}

Theorem \ref{thm:exact} has an exact value with a nonlinear norm term and a nonzero multi-step optimal trajectory. Proposition \ref{prop:conic-program} computes its coefficients from finite conic constraints. Theorem \ref{thm:threshold} identifies the control-bound threshold at which this formula ceases to hold globally.

The two-variable program in \eqref{eq:conic-program} owes its size to orthogonality of $P$ and the shared Euclidean norm in \eqref{eq:dynamics-cost}. The optimal feedback does not saturate the control bound throughout the interior of \eqref{eq:state-cone}. A broader exact formulation would have to handle conic constraints active on an open set of states, or matrices that do not preserve the Euclidean norm. Neither extension is established here. The results establish exactness for the displayed dynamics and stage cost under the stated parameter conditions.

\begin{thebibliography}{9}
\bibitem{pates-rantzer}
R.~Pates and A.~Rantzer,
``Optimal control on positive cones,''
in \emph{Proceedings of the IEEE Conference on Decision and Control}, 2024.
\newblock arXiv:2407.18774.

\bibitem{de-farias-van-roy}
D.~P. de Farias and B.~Van Roy,
``The linear programming approach to approximate dynamic programming,''
\emph{Operations Research}, vol.~51, no.~6, pp.~850--865, 2003.
\end{thebibliography}

\end{document}
